\chapter{Типи нейронів}
\label{sec:neurontypes}

\section{Нейрон як чорна скринька}

Припустімо, вам дали мішок із вісімдесятьма шістьма мільярдами нейронів~\cite{Azevedo2009} і попросили розкласти їх на купки. Як би ви це зробили?

Інженер, якому дали мішок незнайомих мікросхем, не став би сортувати їх за формою корпусу. Він запитав би: скільки в деталі входів, скільки виходів і що вона видає на вихід. Намалюймо нейрон так само --- як блок на схемі (рис.~\ref{fig:blackbox}).

\begin{figure}[t]
\centering
\begin{tikzpicture}[>=Stealth,x=1cm,y=1cm]
% the box
\node[draw,very thick,minimum width=2.6cm,minimum height=2.6cm,fill=blue!6] (N) at (0,0) {};
\node at (0,0.55) {\small нейрон};
\node at (0,-0.15) {$\displaystyle u=\sum_i w_i x_i$};
\node at (0,-0.8) {\small $y=\Theta(u-\theta)$};
% inputs
\foreach \k/\y/\lab in {1/1.05/{x_1},2/0.55/{x_2},3/0.05/{x_3}}{
  \draw[->,thick,blue!60!black] (-4.2,\y) -- (-1.3,\y);
  \node[left] at (-4.2,\y) {$\lab$};
  \node[above,blue!60!black] at (-2.1,\y-0.06) {\scriptsize $w_{\k}$};}
\node at (-4.45,-0.4) {$\vdots$};
\node at (-2.1,-0.4) {$\vdots$};
\draw[->,thick,blue!60!black] (-4.2,-1.05) -- (-1.3,-1.05);
\node[left] at (-4.2,-1.05) {$x_n$};
\node[above,blue!60!black] at (-2.1,-1.11) {\scriptsize $w_{n}$};
\draw[decorate,decoration={brace,amplitude=5pt},gray!60!black] (-4.9,-1.2) -- (-4.9,1.2);
\node[left,align=right,gray!40!black] at (-5.1,0) {\small $n\sim10^3$--$10^5$\\[-2pt]\small входів};
% single output wire, then fan-out
\draw[very thick,red!70!black] (1.3,0) -- (3.2,0);
\foreach \x in {1.6,2.2,2.34,2.9}{\draw[thick,red!70!black] (\x,0) -- (\x,0.4);}
\node[above right,font=\tiny\itshape,gray!30!black,inner sep=1pt] at (1.42,0.45) {клац\dots\ клац-клац\dots};
\node[below,red!70!black,align=center] at (2.25,-0.05) {\scriptsize один вихід $y$};
\fill[red!70!black] (3.2,0) circle (0.06);
\foreach \y in {1.3,0.65,0,-0.65,-1.3}{
  \draw[->,thick,red!70!black] (3.2,0) -- (5.0,\y);}
\draw[decorate,decoration={brace,amplitude=5pt},gray!60!black] (5.3,1.45) -- (5.3,-1.45);
\node[right,align=left,gray!40!black] at (5.5,0) {\small копії $y$\\[-2pt]\small на $10^3$--$10^4$\\[-2pt]\small інших нейронів};
\end{tikzpicture}
\caption{Нейрон очима інженера. Багато входів $x_i$, кожен зі своєю вагою $w_i$; усередині --- зважена сума $u$ і порівняння з порогом $\theta$ ($\Theta$ --- сходинка: $1$, якщо аргумент додатний, і $0$ інакше); один вихід $y$ --- послідовність імпульсів, розмножена розгалуженням на тисячі отримувачів.}
\label{fig:blackbox}
\end{figure}

На виході блока --- не число, а послідовність коротких імпульсів напруги: «клац», пауза, «клац-клац», довга пауза. Усі імпульси однакової висоти, тож уся інформація --- в тому, \emph{коли} вони відбуваються. Усередині блока, в першому грубому наближенні, стоять суматор і компаратор: входи додаються з вагами, і якщо сума перевищила поріг, блок видає імпульс. У наступному розділі ми замінимо його моделлю, яка чесно працює в неперервному часі.

Вихід один, але провід розгалужується, і кожна гілка несе \emph{ту саму} копію сигналу. В електроніці це називають навантажувальною здатністю за виходом, fan-out. У нейрона кори вона становить кілька тисяч; логічний вентиль у мікросхемі зазвичай навантажений на одиниці інших вентилів.

\emph{Що} ж передається вихідним проводом до отримувача? Імпульс добігає до кінця кожної гілки й там перетворюється на хімічний сигнал: гілка викидає у вузьку щілину між клітинами дрібку певної речовини, \emph{нейромедіатора}, яка змінює вхідний струм отримувача. Ось і ознака для сортування: \emph{який медіатор викидає вихід нейрона}.

\section{Один вихід --- один тип сигналу}

Сортування має сенс, лише якщо в кожного нейрона тип виходу один. Чи так це? Дослід каже, що майже так.

\begin{keyblock}
\begin{proposition}[Один нейрон --- один тип виходу]
\label{prop:dale}
Майже кожен нейрон має один головний медіатор і викидає його в усіх гілках свого виходу. Тому тип нейрона задається його головним медіатором~\cite{Strata1999}.
\end{proposition}
\end{keyblock}

Домовмося позначати тип нейрона \emph{першими чотирма літерами англійської назви його головного медіатора}. Так, нейрон, головний медіатор якого глутамат, --- це \nt{GLUT}-нейрон. Усі типи зібрано в таблиці~\ref{tab:types}.

\begin{table}[t]
\centering
\footnotesize
\setlength{\tabcolsep}{3pt}
\renewcommand{\arraystretch}{1.15}
\begin{tabular}{>{\raggedright}p{1.0cm}>{\raggedright}p{2.05cm}>{\raggedright}p{1.75cm}>{\raggedright}p{1.6cm}>{\centering}p{0.7cm}>{\raggedright}p{1.55cm}>{\raggedright}p{1.8cm}>{\raggedright\arraybackslash}p{3.2cm}}
\toprule
Тип & Головний медіатор & Шина & Швидкість & Знак & Нейронів у мозку & Виходів у нейрона & Роль в обчисленні \\
\midrule
\nt{GLUT} & глутамат & адресна & швидко й повільно & $+$ & $\sim8\cdot10^{10}$ & $\sim10^4$ & головна додатна вага \\
\nt{GABA} & ГАМК & адресна & швидко й повільно & $-$ & $\sim3\cdot10^{9}$ & $10^3$--$10^4$ & головна від'ємна вага \\
\nt{GLYC} & гліцин & адресна & швидко & $-$ & $10^6$--$10^7$ & $10^2$--$10^3$ & від'ємна вага в спинному мозку та стовбурі; другий ключ для одного з рецепторів глутамату \\
\nt{ACET} & ацетилхолін & обидві & швидко й повільно & $\pm$ & $\sim10^6$ & м'яз: $10$--$10^3$; мозок: $\sim10^5$ & вихід на м'яз; у мозку --- швидкість навчання (гіпотеза) \\
\nt{DOPA} & дофамін & широко\-мовна & повільно & $\pm$ & $\sim5\cdot10^{5}$ & $10^5$--$10^6$ & помилка передбачення винагороди $\delta$ \\
\nt{SERO} & серотонін & широко\-мовна & повільно (один тип рецептора швидкий) & $\pm$ & $\sim3\cdot10^{5}$ & $\sim10^5$ & горизонт планування (гіпотеза) \\
\nt{HIST} & гістамін & широко\-мовна & повільно & $+$ & $\sim6\cdot10^{4}$ & немає оцінки & глобальний сигнал «не спи» \\
\nt{NORA} & норадреналін & широко\-мовна & повільно & $\pm$ & $\sim5\cdot10^{4}$ & $\sim10^5$ & коефіцієнт підсилення (гіпотеза) \\
\nt{ADRE} & адреналін & широко\-мовна & повільно & $\pm$ & $\sim10^4$ & немає оцінки & мало нейронів; роль у мозку неясна \\
\nt{ASPA} & аспартат & адресна & швидко & $+$ & немає оцінки & немає оцінки & слабка додатна вага; роль спірна \\
\bottomrule
\end{tabular}
\caption{Типи нейронів, за спаданням кількості нейронів у мозку. \emph{Шина}: адресна --- медіатор викидається у вузьку щілину до певного отримувача; широкомовна --- від малої групи нейронів-джерел на великі області (підрозділ~\ref{sec:buses}). \emph{Швидкість}: швидко --- через іонотропні рецептори, мілісекунди; повільно --- через метаботропні, сотні мілісекунд і довше. \emph{Знак} --- звичайний; остаточно його задає отримувач (підрозділ~\ref{sec:receiver}), і $\pm$ означає, що трапляються обидва знаки. \emph{Нейронів у мозку} --- скільки нейронів цього типу в усьому мозку людини, за порядком величини; загалом нейронів близько $8.6\cdot10^{10}$~\cite{Azevedo2009}. Майже всі \nt{GLUT}-нейрони --- близько $7\cdot10^{10}$ --- це дрібні вхідні нейрони мозочка; числа для \nt{GLUT} і \nt{GABA} отримано з цього підрахунку та частки \nt{GABA}-нейронів у корі~\cite{Hendry1987}. З нечисленних типів безпосередньо підраховано лише \nt{HIST}~\cite{Airaksinen1991} і головну з двох груп \nt{DOPA}-нейронів~\cite{Pakkenberg1991}, тож для \nt{DOPA} це нижня межа; для \nt{SERO} --- сума двох часткових підрахунків~\cite{Baker1990,Baker1991}. Числа для \nt{NORA}, \nt{GLYC}, \nt{ACET} і \nt{ADRE} --- грубі оцінки за порядком величини, без прямого підрахунку. \emph{Виходів у нейрона} --- скільки в одного нейрона синаптичних закінчень, тобто місць викиду медіатора на гілках вихідного проводу (рис.~\ref{fig:blackbox}), за порядком величини. Для широкомовних типів закінчення безпосередньо не рахували: для \nt{DOPA} число виведено з того, яку частку своєї мішені покриває розгалуження одного вихідного проводу~\cite{Matsuda2009}, для решти оцінки ще грубіші. Для \nt{ACET} два випадки: нейрон, що керує м'язом, і широкомовний нейрон у мозку.}
\label{tab:types}
\end{table}

% the pie is a table-type float with a figure caption, so it can never float ahead of Table~\ref{tab:types}
\begin{table}[tb]
\makeatletter\def\@captype{figure}\makeatother
\centering
\definecolor{vizglut}{HTML}{2A78D6}
\definecolor{vizgaba}{HTML}{E34948}
\definecolor{vizrest}{HTML}{8A8986}
\begin{tikzpicture}[>=Stealth]
% pie: GLUT 96.4 %, GABA 3.6 % (13.0 degrees), the rest 0.006 % (a hairline at 0 degrees)
\fill[vizglut] (0,0) -- (13:1.9) arc (13:360:1.9) -- cycle;
\fill[vizgaba] (0,0) -- (0:1.9) arc (0:13:1.9) -- cycle;
\draw[white,line width=1pt] (0,0) -- (13:1.9);
\draw[vizrest,line width=1.2pt] (0,0) -- (0:1.9);
\node[white,align=center,font=\small] at (190:0.95) {\nt{GLUT}\\$96.4\,\%$};
\draw[gray!60!black] (6.5:1.75) -- (2.05,2.1);
\node[anchor=west,font=\small] at (2.05,2.1) {\nt{GABA} $3.6\,\%$};
% zoom box for the remaining types
\draw[densely dashed,vizrest] (1.9,0) -- (3.1,1.65);
\draw[densely dashed,vizrest] (1.9,0) -- (3.1,-2.2);
\draw[vizrest,rounded corners=3pt] (3.1,-2.2) rectangle (10.1,1.65);
\node[anchor=west,font=\scriptsize] at (3.2,1.4) {усі інші разом: $\approx0.006\,\%$; логарифмічна шкала};
\foreach \code/\lg/\y/\val/\lx in {GLYC/6.477/1.0/{$10^6$--$10^7$}/4.05,ACET/6.0/0.6/{$\sim10^6$}/3.0,DOPA/5.699/0.2/{$\sim5\cdot10^5$}/2.699,SERO/5.477/-0.2/{$\sim3\cdot10^5$}/2.477,HIST/4.778/-0.6/{$\sim6\cdot10^4$}/1.778,NORA/4.699/-1.0/{$\sim5\cdot10^4$}/1.699,ADRE/4.0/-1.4/{$\sim10^4$}/1.0}{
  \node[anchor=east,font=\scriptsize] at (4.35,\y) {\nt{\code}};
  \fill[vizrest] (4.45,\y-0.13) rectangle ({4.45+\lg-3},\y+0.13);
  \node[anchor=west,font=\scriptsize] at ({4.5+\lx},\y) {\val};}
\draw[vizrest,thick] (7.45,1.0) -- (8.45,1.0);
\draw[vizrest,thick] (7.45,0.92) -- (7.45,1.08);
\draw[vizrest,thick] (8.45,0.92) -- (8.45,1.08);
\draw[gray!60!black] (4.45,-1.72) -- (8.45,-1.72);
\foreach \e in {3,4,5,6,7}{
  \draw[gray!60!black] ({4.45+\e-3},-1.72) -- ({4.45+\e-3},-1.66);
  \node[below,font=\scriptsize] at ({4.45+\e-3},-1.72) {$10^{\e}$};}
\end{tikzpicture}
\caption{Частки типів нейронів у мозку за оцінками таблиці~\ref{tab:types}. Коло майже цілком зайняте \nt{GLUT}; \nt{GABA} --- вузький сектор, а всі інші типи разом --- волосина на межі. Праворуч цю волосину збільшено, за логарифмічною шкалою кількості нейронів; для \nt{GLYC} стовпчик сягає середини діапазону $10^6$--$10^7$, відрізок показує весь діапазон. Для \nt{ASPA} оцінки немає, і на рисунку його немає.}
\label{fig:typepie}
\end{table}

Наскільки нерівні ці купки, видно на рис.~\ref{fig:typepie}: з кожної тисячі нейронів $964$ --- \nt{GLUT}, $36$ --- \nt{GABA}, а на всі інші типи разом припадає близько шести нейронів на сто тисяч.

Назва GABA (англійське скорочення ГАМК) і так складається з чотирьох літер. Речовини, які майже ніколи не бувають головним медіатором, --- нейропептиди, гази, ліпіди, пурини, --- власного коду не отримують; про них --- у додатку~\ref{sec:othermediators}. Слово «майже» у твердженні~\ref{prop:dale} важливе. Багато нейронів викидають разом із головним медіатором допоміжні речовини~\cite{Yao2023}. Деякі викидають і другий швидкий медіатор, навіть протилежного знака: частина \nt{DOPA}-нейронів викидає ще й ГАМК~\cite{Tritsch2012}. А в деяких нейронів різні медіатори виходять із різних гілок одного й того самого виходу~\cite{Zhang2015}. Усе це виміряно, тож «один нейрон --- один медіатор» --- правило з винятками, а не закон. Але як перший поділ воно витримує перевірку: у клітинному атласі мозку миші, де нейрони розкладено за генами, що в них працюють, на понад п'ять тисяч груп, нейрони так само діляться на великі класи насамперед за головним медіатором~\cite{Yao2023}.

Із цього правила випливає наслідок, який одразу відрізняє мозок від штучної нейронної мережі. У штучній мережі той самий нейрон може мати додатну вагу до одного отримувача й від'ємну до іншого: ваги $w_{ij}$ --- просто числа в матриці. У мозку знак дії здебільшого визначається медіатором, а медіатор у нейрона один. Отже, якщо нейрон $j$ збуджує одного отримувача, він збуджує й усіх інших. Увесь \emph{стовпець} матриці ваг, що виходить із нейрона $j$, одного знака:
\begin{empheq}[box=\keybox]{equation}
\operatorname{sign} w_{ij} \;=\; s_j \quad\text{для всіх отримувачів } i, \qquad s_j\in\{+1,-1\}.
\label{eq:dalesign}
\end{empheq}
Нейрони поділяються на збуджувальні ($s_j=+1$) і гальмівні ($s_j=-1$), і це властивість самого нейрона, а не зв'язку. Якщо мережі потрібен від'ємний зв'язок від збуджувального нейрона, його доводиться робити через посередника: збуджувальний нейрон збуджує гальмівний, а той гальмує ціль: від'ємна вага із затримкою на одну ланку.

Медіаторів відомо понад сотню, але майже вся робота мозку тримається на півдюжині, і вони поділяються на два зовсім різні класи.

\section{Дві шини: адресна й широкомовна}
\label{sec:buses}

У комп'ютерних мережах є два способи доставити повідомлення. Перший --- \emph{адресний}, unicast: пакет іде від одного відправника до певного отримувача, швидко, і більше його ніхто не бачить. Другий --- \emph{широкомовний}, broadcast: відправник передає одне повідомлення відразу всім вузлам мережі. Нейрони користуються обома способами, і медіатори поділяються саме за цією ознакою (рис.~\ref{fig:phone_radio}).

\begin{figure}[t]
\centering
\begin{tikzpicture}[>=Stealth,scale=0.95]
% ---- unicast: point-to-point wiring
\node[above] at (2.0,3.3) {\small \textbf{Адресна}: глутамат і ГАМК};
\foreach \i in {0,...,4}{
  \fill[blue!60!black] (0.2+0.9*\i,2.5) circle (0.1);
  \fill[blue!60!black] (0.2+0.9*\i,0.6) circle (0.1);}
\foreach \a/\b in {0/0,0/1,1/1,1/3,2/2,3/2,3/4,4/4,2/0}{
  \draw[->,blue!60!black] (0.2+0.9*\a,2.38) -- (0.2+0.9*\b,0.72);}
\fill[red!70!black] (1.1,1.55) circle (0.09);
\pic[scale=1.2] at (1.75,1.9) {letter};
\draw[->,red!70!black,thick] (1.1,1.55) -- (0.28,0.7);
\draw[->,red!70!black,thick] (1.1,1.55) -- (1.95,0.72);
\node[align=center,below] at (2.0,0.2) {\small мільярди нейронів,\\[-2pt]\small у кожного свої адресати,\\[-2pt]\small час: мілісекунди};
% ---- broadcast: one small source, global targets
\begin{scope}[xshift=7.2cm]
\node[above] at (1.8,3.3) {\small \textbf{Широкомовна}: дофамін, \dots};
\foreach \i in {0,...,8}{\fill[blue!60!black] (-0.2+0.5*\i,2.75) circle (0.08);}
\fill[orange!85!black] (1.8,0.35) ellipse (0.35 and 0.18);
\pic[scale=1.5] at (2.7,0.75) {mast};
\foreach \x in {-0.2,0.3,0.8,1.3,1.8,2.3,2.8,3.3,3.8}{
  \draw[->,orange!85!black] (1.8,0.5) .. controls (1.8,1.3) and (\x,1.6) .. (\x,2.62);}
\node[align=center,below] at (1.8,0.1) {\small тисячі --- сотні тисяч нейронів,\\[-2pt]\small одне повідомлення всім,\\[-2pt]\small час: сотні мілісекунд і довше};
\end{scope}
\end{tikzpicture}
\caption{Два способи передачі. Ліворуч --- адресна шина, схожа на пошту з адресою на конверті; червоним виділено один нейрон і двох його отримувачів. Праворуч --- широкомовна, як радіостанція: невелика група нейронів-джерел, і те саме повідомлення отримують усі.}
\label{fig:phone_radio}
\end{figure}

\emph{Адресна шина} --- це глутамат і ГАМК. Їх викидає переважна більшість нейронів. Кожен вихід веде до певних клітин, медіатор діє лише у вузькій щілині контакту, і діє швидко: за мілісекунду відкриваються іонні канали отримувача, за кілька мілісекунд усе закінчується. Це шина \emph{даних}: нею передається те, що мозок обчислює.

\emph{Широкомовна шина} --- дофамін, серотонін, норадреналін і почасти ацетилхолін. Їх викидають крихітні групи нейронів, але вихід кожного з них розгалужується неймовірно широко. Медіатор часто викидається не у вузьку щілину, а в міжклітинний простір, де він розпливається й діє на всіх, хто поруч. Діє повільно: не відкриває канали безпосередньо, а запускає всередині отримувача ланцюжок хімічних реакцій. Нею йдуть не дані, а \emph{керувальні параметри}, спільні для великих ділянок мережі, які змінюють режим її роботи: коефіцієнт підсилення, швидкість навчання, сигнал «це було добре». Тому такі медіатори називають \emph{модуляторами}. Довго вважали, що кожен такий канал --- одне число на весь мозок. Сучасні вимірювання уточнили картину: канал --- не один провід, а невеликий джгут паралельних ліній. Нейрони-джерела одного медіатора поділяються на підсистеми, що мають своїх отримувачів і своє повідомлення, а скільки медіатора викинути на місці, підлаштовують ще й місцеві сигнали~\cite{Ren2018,Poe2020}. Таких ліній --- одиниці чи десятки, а не мільярди, тож канал лишається широкомовним.

Якщо вам потрібна одна картинка з цього розділу --- ось вона. Мозок --- величезна мережа з адресною передачею даних, над якою висить кілька широкомовних каналів із керувальними сигналами. Програміст упізнає знайомий устрій: обчислення плюс невеликий набір керувальних змінних, кожна з яких спільна для великої ділянки мережі.

\section{\nt{GLUT}: додатна вага}

Глутамат --- головний збуджувальний медіатор мозку. Коли вихід нейрона викидає глутамат, в отримувача відкриваються канали, через які всередину входить натрій, напруга на мембрані отримувача зростає, і його шанс видати власний імпульс збільшується. Мовою рис.~\ref{fig:blackbox} це вхід із \emph{додатною} вагою, $w_i>0$.

Хто видає глутамат? Майже всі, хто передає дані на великі відстані: від трьох чвертей до чотирьох п'ятих нейронів кори й більшість нейронів, що з'єднують різні області мозку. Мережа мозку --- це насамперед мережа глутаматних зв'язків.

Одна інженерна деталь. Після викиду концентрація глутамату в щілині злітає в тисячі разів, приблизно до мілімоля, і спадає зі сталою часу близько мілісекунди~\cite{Clements1992}: глутамат розпливається з щілини, а спеціальні білки-насоси відкачують його назад у клітини. Навіщо? Із тієї самої причини, з якої в лінії зв'язку після кожного символу сигнал повертають до нуля. Якщо медіатор не прибрати, наступний імпульс прийде на «зайняту» лінію, й імпульси з інтервалом у кілька мілісекунд зіллються.

\section{\nt{GABA}: від'ємна вага}

Уявіть мережу, у якій усі ваги додатні. Якщо в середньому кожен імпульс породжує більше ніж один наступний імпульс, активність зростає експоненційно й за кілька кроків охоплює всю мережу. Електронник упізнає петлю додатного зворотного зв'язку з підсиленням більше одиниці: схема йде в насичення. Щоб цього не сталося, потрібні від'ємні ваги. Їхнє головне джерело --- гамма-аміномасляна кислота, скорочено ГАМК.

ГАМК відкриває в отримувача канали, що пропускають іони хлору. Рівноважний потенціал хлору лежить біля потенціалу спокою або трохи нижче, тому відкриті хлорні канали утримують мембрану біля спокою й не дають підняти напругу до порога. Це вхід із \emph{від'ємною} вагою, $w_i<0$. ГАМКергічні --- від п'ятої частини до чверті нейронів кори~\cite{Hendry1987}. Їхні виходи короткі, і гальмують вони своїх найближчих сусідів.

Гальмування в мозку робить набагато більше, ніж просто віднімання. Воно працює як від'ємний зворотний зв'язок, що тримає всю мережу в робочій точці. Вважають, що в нормально працюючій корі збуджувальний і гальмівний струми, які надходять на нейрон, майже точно врівноважують один одного: такий режим передбачено теорією~\cite{vanVreeswijk1996} і видно у вимірюваннях струмів у живій корі~\cite{Haider2006}, і нейрон весь час сидить біля порога. Схема, стабілізована сильним від'ємним зворотним зв'язком біля нестійкої точки, швидко й чутливо реагує на малі сигнали --- давній інженерний прийом.

\section{\nt{DOPA}: сигнал помилки}

Перший широкомовний канал --- дофамін. Дофамінових нейронів у людини близько пів мільйона, приблизно один на сто сімдесят тисяч; стільки підраховано в головній із двох їхніх груп~\cite{Pakkenberg1991}, тож це нижня межа. Їхні виходи розходяться до тих областей, які обирають дії.

Що передає цей канал? Відповідь дають досліди, у яких реєструють імпульси дофамінових нейронів у тварини, що отримує винагороду~\cite{Schultz1997}. Тварині час від часу несподівано дають краплю соку, і дофамінові нейрони відповідають на неї короткою пачкою імпульсів. Потім перед соком починають запалювати лампочку, завжди за секунду до соку. Коли тварина вивчає цей зв'язок, нейрони починають відповідати пачкою на \emph{лампочку}, а на сам сік, що прийшов точно вчасно, не відповідають зовсім. А якщо після лампочки соку не дають, то в момент, коли він мав прийти, нейрони \emph{замовкають}, і їхня частота падає нижче звичайної.

Правило, що пояснює всі три випадки (рис.~\ref{fig:rpe}): нейрони повідомляють не про те, наскільки добре, а про те, наскільки \emph{краще чи гірше, ніж очікувалося}.

\begin{figure}[t]
\centering
\begin{tikzpicture}[>=Stealth,x=3.4cm,y=1cm]
\foreach \y/\lab in {2.8/{несподіваний сік},1.4/{сік після лампочки},0/{лампочка, а соку немає}}{
  \draw[gray!70!black] (0,\y) -- (3,\y);
  \draw[densely dotted,gray!60!black] (0,\y+0.3) -- (3,\y+0.3);
  \node[left,align=right,font=\small] at (-0.03,\y+0.35) {\lab};}
% row 1: juice without a cue -> burst at the juice
\draw[very thick,orange!85!black] plot[domain=0:3,samples=120] (\x,{2.8+0.3+0.75*exp(-((\x-2.08)/0.07)^2)});
\pic[scale=0.8] at (2,2.58) {juice};
% row 2: cue then juice -> burst at the cue, nothing at the juice
\draw[very thick,orange!85!black] plot[domain=0:3,samples=120] (\x,{1.4+0.3+0.75*exp(-((\x-1.08)/0.07)^2)});
\pic[scale=0.85] at (1,1.15) {lamp}; \pic[scale=0.85] at (2,1.15) {juice};
% row 3: cue, juice omitted -> burst at the cue, dip when the juice was due
\draw[very thick,orange!85!black] plot[domain=0:3,samples=120] (\x,{0.3+0.75*exp(-((\x-1.08)/0.07)^2)-0.28*exp(-((\x-2.1)/0.12)^2)});
\pic[scale=0.85] at (1,-0.25) {lamp}; \pic[scale=0.85] at (2,-0.25) {nojuice};
\node[right,font=\scriptsize] at (2.36,0.1) {провал};
\node[right,font=\scriptsize] at (2.13,3.75) {сюрприз!};
\node[left,font=\scriptsize] at (0.97,2.25) {сплеск на лампочку};
\node[above,font=\scriptsize] at (2,1.75) {нічого};
\draw[->,gray!70!black] (0,-0.75) -- (3.05,-0.75) node[right,black,font=\small] {$t$};
\draw[gray!60!black] (1,-0.8) -- (1,-0.7) node[below=2pt,font=\scriptsize] {лампочка, $1\,\text{с}$};
\draw[gray!60!black] (2,-0.8) -- (2,-0.7) node[below=2pt,font=\scriptsize] {сік, $2\,\text{с}$};
\end{tikzpicture}
\caption{Частота імпульсів \nt{DOPA}-нейронів у трьох дослідах (схематично, за~\cite{Schultz1997}). Пунктир --- рівна фонова частота. Лампочка --- сигнал, крапля --- сік, перекреслена крапля --- сік, який обіцяли, але не дали. Відгук --- помилка передбачення~\eqref{eq:rpe}.}
\label{fig:rpe}
\end{figure}

Програміст, знайомий із навчанням із підкріпленням~\cite{SuttonBarto2018}, упізнає цю величину відразу. Агент, який вчиться обирати дії, зберігає оцінку $V(s)$ --- скільки винагороди він очікує, перебуваючи в стані $s$. Після кожного кроку він порівнює очікування з тим, що отримав:
\begin{empheq}[box=\keybox]{equation}
\delta_t \;=\; r_t \;+\; \gamma\,V(s_{t+1}) \;-\; V(s_t) ,
\label{eq:rpe}
\end{empheq}
і виправляє оцінку: $V(s_t)\leftarrow V(s_t)+\alpha\,\delta_t$. Тут $r_t$ --- отримана винагорода, $\gamma<1$ --- коефіцієнт дисконтування (наскільки майбутня винагорода дешевша за негайну), $\alpha$ --- швидкість навчання. Величину $\delta_t$ називають \emph{помилкою часових різниць}.

Вона поводиться точнісінько як дофамінові нейрони. Несподіваний сік: $V$ ще нуль, $r>0$, $\delta>0$. Вивчений сік: винагороду передбачила лампочка, $V(s_t)$ перед соком уже дорівнює $r$, і $\delta=0$. Сік не прийшов: $V(s_t)=r$, а $r_t=0$, $\delta<0$. А сплеск переїжджає на лампочку, бо саме в момент лампочки очікування стрибком зросло: $\gamma V(s_{t+1})-V(s_t)>0$. Кроки $t,t+1,\dots$ тут --- умовність алгоритму: в організмі немає годинника, що відлічує кроки, і ту саму величину в неперервному часі ми отримаємо в розділі~\ref{sec:dofa}.

Отже, дофамін --- це широкомовна розсилка скалярного сигналу помилки $\delta$ усім, хто має за ним навчатися. У першому наближенні --- один провід на весь мозок і одне число в кожен момент часу; наскільки цей провід насправді джгут, розглянуто в розділі~\ref{sec:dofaalt}.

\section{Інші широкомовні канали}

Для решти модуляторів є лише робочі гіпотези, які перекладають їхні ролі на ту саму мову глобальних параметрів~\cite{Doya2002}. Сприймайте їх саме як гіпотези.

\emph{\nt{NORA}, норадреналін: коефіцієнт підсилення.} Норадреналінових нейронів ще менше, ніж дофамінових, за грубою оцінкою кілька десятків тисяч, а їхні виходи розходяться практично по всьому мозку. Вони різко активуються, коли стається щось несподіване або важливе. Норадреналін змінює крутизну передавальної характеристики нейронів-отримувачів: слабкі входи стають слабшими, сильні --- сильнішими. Вимірюють це так: фонові імпульси отримувача пригнічуються сильніше, ніж відповідь на вхід, і відношення сигналу до шуму зростає~\cite{Foote1975NE}. У простій моделі це описують одним числом: якщо нейрон відповідає частотою $f=F\bigl(g\cdot u\bigr)$ на вхідний струм $u$, то норадреналін збільшує коефіцієнт $g$~\cite{AstonJones2005} (докладніше --- у розділі~\ref{sec:nora}). Мережа з великим $g$ працює «контрастніше» й швидше ухвалює рішення; з малим --- «розмито» й більше перебирає варіанти, як агент навчання з підкріпленням у режимі пошуку.

\emph{\nt{ACET}, ацетилхолін: швидкість навчання.} Ацетилхолін керує тим, наскільки мережа слухає поточний вхід, а наскільки --- власні збережені дані. За високого рівня ацетилхоліну входи від органів чуття посилюються, а внутрішні, «пригадувальні» зв'язки слабшають, і водночас полегшується зміна ваг~\cite{Hasselmo2006}. Грубо кажучи, це перемикач «запис/читання» і разом із ним --- швидкість навчання $\alpha$.

\emph{\nt{SERO}, серотонін: горизонт планування.} Що передає серотонін, зрозуміло найгірше. Одна з гіпотез пов'язує його з коефіцієнтом дисконтування $\gamma$ з~\eqref{eq:rpe}: високий рівень серотоніну відповідає готовності почекати на велику винагороду, а не хапати малу зараз. Серотонінові нейрони працюють рівно й повільно, кілька імпульсів за секунду в неспанні, і майже замовкають в одній із фаз сну~\cite{JacobsAzmitia1992}.

Усі шість медіаторів зібрано в таблиці~\ref{tab:transmitters}.

\begin{table}[t]
\centering
\small
\begin{tabular}{>{\raggedright}p{2.4cm}>{\raggedright}p{3.0cm}>{\raggedright}p{2.0cm}>{\raggedright\arraybackslash}p{5.2cm}}
\toprule
Тип нейрона & Частка нейронів & Шина & Роль в обчисленні \\
\midrule
\nt{GLUT} (глутамат) & більшість & адресна & додатна вага, $w>0$ \\
\nt{GABA} (ГАМК) & 20--25\% у корі & адресна & від'ємна вага, $w<0$; стабілізація \\
\nt{DOPA} (дофамін) & $\sim 6\cdot10^{-6}$ & широко\-мовна & помилка передбачення $\delta$ \\
\nt{NORA} (норадреналін) & $\sim 6\cdot10^{-7}$ & широко\-мовна & коефіцієнт підсилення $g$ (гіпотеза) \\
\nt{ACET} (ацетилхолін) & мала & обидві & швидкість навчання $\alpha$; «запис/читання» (гіпотеза) \\
\nt{SERO} (серотонін) & $\sim 3\cdot10^{-6}$ & широко\-мовна & дисконтування $\gamma$ (гіпотеза) \\
\bottomrule
\end{tabular}
\caption{Головні медіатори мозку мовою обчислень. Правий стовпець --- сильне спрощення: він надійний для глутамату, ГАМК і дофаміну, для решти це гіпотези.}
\label{tab:transmitters}
\end{table}

\section{Знак задає отримувач}
\label{sec:receiver}

Я трохи вас обдурив, коли казав: глутамат --- додатна вага, ГАМК --- від'ємна. Жодна молекула не несе в собі знака. Молекула --- це просто ключ. Що станеться, коли її впіймають, вирішує \emph{замок} --- рецептор на клітині-отримувачі.

Найкращий приклад --- дофамін. Він має рецептори двох родин~\cite{Beaulieu2011}. На одних він полегшує збудження клітини, на інших --- утруднює. В області, яка обирає дії, одні нейрони несуть рецептори першого типу, інші --- другого, і той самий дофаміновий сигнал водночас підсилює одну групу й послаблює іншу. Це як один широкомовний пакет, який різні вузли мережі інтерпретують по-різному.

\begin{keyblock}
\begin{proposition}[Зміст повідомлення задає отримувач]
\label{prop:receptor}
Знак і швидкість дії медіатора визначає не сам медіатор, а рецептор, з яким він зв'язується. Рецептори бувають двох родів. \emph{Іонотропні} самі є іонними каналами й відкриваються за частки мілісекунди: це пряме керування струмом, як затвор транзистора. \emph{Метаботропні} запускають усередині клітини ланцюжок хімічних реакцій і діють за сотні мілісекунд і довше: це радше запис у регістр налаштувань, який змінює подальшу роботу клітини. Адресна шина говорить здебільшого через перші, широкомовна --- здебільшого через другі.
\end{proposition}
\end{keyblock}

Чому ж тоді взагалі можна сказати «глутамат збуджує»? Тому що майже всі рецептори глутамату, які трапляються на практиці, збуджувальні, а майже всі рецептори ГАМК --- гальмівні. Це не закон природи, а дуже надійна домовленість, і правило~\eqref{eq:dalesign} на ній і тримається.

\section{Що взяти з собою}

Переважна більшість нейронів --- адресна шина даних із вагами одного знака на нейрон; мізерна меншість --- широкомовні канали, що передають на великі ділянки мозку кілька спільних керувальних сигналів. Близько двох нейронів із кожних ста тисяч --- і від них залежить, як навчається і в якому режимі працює вся решта мережі. Для інженера це не дивно: у будь-якому комп'ютері регістрів керування незрівнянно менше, ніж комірок пам'яті, і все ж без них машина не працює.

У наступному розділі ми перевіримо, що може обчислити мережа з самих лише \nt{GLUT}-нейронів, найчисленнішого типу, яка працює в неперервному часі.

\section{Задачі}

\begin{exercise}[\lvlA]
\label{ex:01:1}
\emph{Купки й проводи.} За таблицею~\ref{tab:types} (середина діапазону за логарифмічною шкалою; \nt{ACET} --- $10^5$ виходів): (а)~якщо нейрон --- людина, скільки населень Землі становитимуть \nt{GLUT}-нейрони і місто якого розміру --- усі широкомовні? (б)~У скільки разів частка закінчень \nt{DOPA}, \nt{SERO}, \nt{NORA}, \nt{ACET} більша за частку їхніх нейронів?
\end{exercise}

\begin{exercise}[\lvlA]
\label{ex:01:2}
\emph{Повернення до нуля.} Глутамат у щілині спадає як $c_0e^{-t/\tau_c}$, $c_0=1\mM$, $\tau_c=1\ms$; отримувач помічає концентрацію вище $10$~мкМ. (а)~Скільки часу лінія «зайнята» після одного викиду? (б)~Насоси частково заблоковано, $\tau_c=10\ms$. До якої частоти викидів лінія ще встигає звільнитися?
\end{exercise}

\begin{exercise}[\lvlB]
\label{ex:01:3}
\emph{Куди переїжджає сплеск.} У досліді рис.~\ref{fig:rpe} після лампочки йдуть стани $s_0\to\dots\to s_{K-1}$ із кроком $\Delta$, $T=K\Delta$; винагорода $r$ приходить на останньому переході, момент лампочки непередбачуваний. Знайдіть вивчені $V(s_k)$ і відповідь $\delta$ на лампочку. Поклавши $\gamma=e^{-\Delta/\tau_\gamma}$, знайдіть $\tau_\gamma$ при $T=1$~с, $\Delta=0.1$~с, $\gamma=0.95$.
\end{exercise}

\begin{exercise}[\lvlB]
\label{ex:01:4}
\emph{Моделювання: навчання за помилкою.} Змоделюйте ланцюжок задачі~\ref{ex:01:3} при $K=10$, $\gamma=0.95$, $r=1$ із правилом $V(s_t)\leftarrow V(s_t)+\alpha\delta_t$. У якому випробуванні $n^*$ відповідь на лампочку вперше переганяє відповідь на сік при $\alpha=0.05$, $0.1$, $0.2$, $0.5$? Як $n^*$ залежить від $\alpha$?
\end{exercise}

\begin{exercise}[\lvlB]
\label{ex:01:5}
\emph{Проєктування: розсилка одного числа.} Число $\delta$ треба доставити всім $10^{10}$ нейронам; кожен отримувач, включно з ретрансляторами, має отримати $10$ закінчень від попереднього шару, ланка додає $1\ms$. Скільки нейронів і ланок у найощаднішому дереві ретрансляторів при $10^4$ виходах на нейрон (\nt{GLUT}) і при $10^{5.5}$ (\nt{DOPA})?
\end{exercise}

\begin{exercise}[\lvlB]
\label{ex:01:6}
\emph{Чому рівновага чутлива.} У мережі $80\,\%$ \nt{GLUT} і $20\,\%$ \nt{GABA}, кожен зв'язаний із кожним вагами $J_E/N$ і $-J_I/N$; $\tau\dot r=-r+Wr+h$. При $J_E=10$ єдине ненульове власне значення $W$ дорівнює $0.5$. Знайдіть відношення гальмівного струму до збуджувального і на скільки відсотків досить послабити гальмування, щоб мережа пішла в лавину.
\end{exercise}

\begin{exercise}[\lvlC]
\label{ex:01:7}
\emph{Дослідження: \nt{SERO} --- це $\gamma$?} За задачею~\ref{ex:01:3} відповідь \nt{DOPA}-нейронів на лампочку дорівнює $re^{-T/\tau_\gamma}$. Затримки $T=1,\dots,5$~с по $n$ пред'явлень, $\ln\delta$ вимірюється з розкидом $0.5$. Скільки потрібно $n$, щоб відрізнити $\tau_\gamma=2$~с від $2.4$~с (рівень $0.05$, потужність $0.8$)? Який механізм, крім $\gamma$, дав би такий самий спад із $T$?
\end{exercise}
